
\subsection{P2P, DHT, Chord}

Lorsque l'index prend encore trop de place, il est utile de le r�partir sur un ensemble de serveur qui r�partiront ainsi la charge. Chaque terme, associ� � sa \textit{posting list} est envoy� sur un serveur. Pour cela, nous utiliserons un r�seau P2P Chord avec une DHT.

Notre DHT utilisera comme fonction de hachage, sur le $docId$, \textbf{modulo (N)}. N �tant le nombre de noeuds du r�seau pair-�-pair. Chaque noeud reli� � deux autres noeuds, formant un anneau complet. Afin de parcourir l'anneau, chaque noeud $N_i$ a une partie de la table de hachage globale (d'o� le $D$istributed $H$ash $T$able). La table de $N_i$ contient un certain nombre de pointeurs vers des noeuds suivants dont les identifiants sont :\\
 $N_i$+2$^0$, $N_i$+2$^1$, $N_i$+2$^0$, $N_i$+2$^2$, $N_i$+2$^3$, $N_i$+2$^4$, $N_i$+2$^5$.

Ainsi, lorqu'� un noeud  $N_i$ on cherche une cl� $j$ (recherche ou insertion), si $j$ est sup�rieure � $N_i$+2$^5$ on sautera les 2$^5$ noeuds suivants. S'il est inf�rieur, on test $N_i$+2$^4$. Et ceci r�cursivement jusqu'� atteindre le noeud $N_j$.

\begin{enumerate}
	\item Construire un r�seau Chord de 18 noeuds.
	\item Placer les \textit{Postings List} de la section~\ref{LI}, question~\ref{LI_q} dans le r�seau. Le $termId$ de chaque terme correspond � son ordre d'apparition.\\
France : 1, R�cession : 2, Ralentie : 3, INSEE : 4, USA : 5, Finance : 6, Paulson : 7, Crise : 8, Gouvernement : 9, PIB : 10
	\item Le noeud 1 souhaite ajouter le document 5 contenant les mots suivants (R�cession, Gouvernement). D�crire le d�roulement de l'insertion.
	\item Une requ�te est pos�e au niveau du noeud 14. D�crire les �tapes de l'�valuation de la requ�te suivante : $crise \land ralentie$.
\end{enumerate}